% --- Archon physics formalization source begin ---
% archon:physics
% archon:covers PhyXMiniProblems/problem_phyx_mini_0283.lean
% archon:source-report reports/phyx_mini/problem_phyx_mini_0283.source.json

\chapter{Physics problem phyx\_mini\_0283}
\label{ch:PhyXMiniProblems_problem_phyx_mini_0283}

\paragraph{Problem source.}
\#\# Physical scenario

A common device for entertaining a toddler is a jump seat that hangs from the horizontal portion of a doorframe via elastic cords. Assume that only one cord is on each side in spite of the more realistic arrangement shown. When a child is placed in the seat, they both descend by a distance \$d\_s\$ as the cords stretch (treat them as springs). Then the seat is pulled down an extra distance \$d\_m\$ and released, so that the child oscillates vertically, like a block on the end of a spring. Suppose you are the safety engineer for the manufacturer of the seat. You do not want the magnitude of the child's acceleration to exceed \$0.20g\$ for fear of hurting the child's neck.

\#\# Question

If \$d\_m = 10\$ cm, what value of \$d\_s\$ corresponds to that acceleration magnitude?

\#\# Answer choices

- A: A: 40 cm

- B: B: 45 cm

- C: C: 55 cm

- D: D: 50 cm

\#\# Figure caption (auxiliary; use the image as primary evidence)

The image depicts a child seated in a swing. The swing is suspended from the ceiling by two ropes. The child is wearing a purple top and appears to be sitting comfortably with their feet visible. The background shows a simple ceiling and floor, indicating an indoor setting. The colors are soft, mainly blue and purple tones, and the style is illustrative. There is no text present.

\#\# Dataset metadata

Resonance and Harmonics; Physical Model Grounding Reasoning, Multi-Formula Reasoning

\paragraph{Recorded answer/context.}
D: D: 50 cm

\paragraph{Figure/image path.}
/root/proposal\_for\_physic/science-mango/phyx\_mini\_run/phyx\_data/test\_image/283.png

\paragraph{Formalization target.}
create a compiling Lean file with sorry bodies at `PhyXMiniProblems/problem\_phyx\_mini\_0283.lean`. The Lean declarations must preserve the physical quantities, dimensions or dimensional roles, figure labels, governing-law hypotheses, and final relation expressed by this problem.
Use Mathlib/Physlib names found through LeanExplore where available. If a physics API is missing, introduce faithful local abstractions rather than scalar placeholder aliases.

\begin{theorem}[Physics formalization target]
\label{thm:physics:phyx_mini_0283:target}
\lean{PhyXMiniProblems.ProblemPhyXMini0283.problem_phyx_mini_0283}
\uses{def:physics:phyx-mini-0283:phyxminiproblems-problemphyxmini0283-lengthincentimeters, def:physics:phyx-mini-0283:phyxminiproblems-problemphyxmini0283-jumpseatoscillatorsetup, def:physics:phyx-mini-0283:phyxminiproblems-problemphyxmini0283-matchesproblemandprimaryfigure, def:physics:phyx-mini-0283:phyxminiproblems-problemphyxmini0283-hasphysicaljumpseatparameters, def:physics:phyx-mini-0283:phyxminiproblems-problemphyxmini0283-satisfiestwocordjumpseatlaws, def:physics:phyx-mini-0283:phyxminiproblems-problemphyxmini0283-recordedanswerchoice, def:physics:phyx-mini-0283:phyxminiproblems-problemphyxmini0283-matchesanswerchoice}
The assigned autoformalize agent should translate this physics problem into Lean declarations in the covered file, with theorem and lemma proof bodies written as `by sorry`.
\end{theorem}
\begin{proof}
This is an autoformalization task, not a proof task. Produce statements that can later be proved by the physics prover without weakening the physical model.
\end{proof}

% --- Archon declaration topology begin ---
\section{Declaration topology}
\begin{definition}[Mass Quantity]
  \label{def:physics:phyx-mini-0283:phyxminiproblems-problemphyxmini0283-massquantity}
  \lean{PhyXMiniProblems.ProblemPhyXMini0283.MassQuantity}
  A nonnegative physical mass
\end{definition}
\begin{proof}
  This is the declared model definition of Mass Quantity.
\end{proof}
\begin{definition}[Length Quantity]
  \label{def:physics:phyx-mini-0283:phyxminiproblems-problemphyxmini0283-lengthquantity}
  \lean{PhyXMiniProblems.ProblemPhyXMini0283.LengthQuantity}
  A nonnegative physical length
\end{definition}
\begin{proof}
  This is the declared model definition of Length Quantity.
\end{proof}
\begin{definition}[Acceleration Magnitude Quantity]
  \label{def:physics:phyx-mini-0283:phyxminiproblems-problemphyxmini0283-accelerationmagnitudequantity}
  \lean{PhyXMiniProblems.ProblemPhyXMini0283.AccelerationMagnitudeQuantity}
  A nonnegative acceleration magnitude, with dimension L T⁻²
\end{definition}
\begin{proof}
  This is the declared model definition of Acceleration Magnitude Quantity.
\end{proof}
\begin{definition}[Spring Stiffness Quantity]
  \label{def:physics:phyx-mini-0283:phyxminiproblems-problemphyxmini0283-springstiffnessquantity}
  \lean{PhyXMiniProblems.ProblemPhyXMini0283.SpringStiffnessQuantity}
  A nonnegative spring stiffness, with dimension force per length, M T⁻²
\end{definition}
\begin{proof}
  This is the declared model definition of Spring Stiffness Quantity.
\end{proof}
\begin{definition}[quantity Readout]
  \label{def:physics:phyx-mini-0283:phyxminiproblems-problemphyxmini0283-quantityreadout}
  \lean{PhyXMiniProblems.ProblemPhyXMini0283.quantityReadout}
  Read a nonnegative dimensionful quantity in a coherent unit system
\end{definition}
\begin{proof}
  This is the declared model definition of quantity Readout.
\end{proof}
\begin{definition}[mass In Kilograms]
  \label{def:physics:phyx-mini-0283:phyxminiproblems-problemphyxmini0283-massinkilograms}
  \lean{PhyXMiniProblems.ProblemPhyXMini0283.massInKilograms}
  \uses{def:physics:phyx-mini-0283:phyxminiproblems-problemphyxmini0283-massquantity, def:physics:phyx-mini-0283:phyxminiproblems-problemphyxmini0283-quantityreadout}
  Kilogram readout of a physical mass
\end{definition}
\begin{proof}
  This is the declared model definition of mass In Kilograms.
\end{proof}
\begin{definition}[length In Meters]
  \label{def:physics:phyx-mini-0283:phyxminiproblems-problemphyxmini0283-lengthinmeters}
  \lean{PhyXMiniProblems.ProblemPhyXMini0283.lengthInMeters}
  \uses{def:physics:phyx-mini-0283:phyxminiproblems-problemphyxmini0283-lengthquantity, def:physics:phyx-mini-0283:phyxminiproblems-problemphyxmini0283-quantityreadout}
  Meter readout of a physical length
\end{definition}
\begin{proof}
  This is the declared model definition of length In Meters.
\end{proof}
\begin{definition}[centimeter Unit Choices]
  \label{def:physics:phyx-mini-0283:phyxminiproblems-problemphyxmini0283-centimeterunitchoices}
  \lean{PhyXMiniProblems.ProblemPhyXMini0283.centimeterUnitChoices}
  SI base units with centimeters selected as the length unit
\end{definition}
\begin{proof}
  This is the declared model definition of centimeter Unit Choices.
\end{proof}
\begin{definition}[length In Centimeters]
  \label{def:physics:phyx-mini-0283:phyxminiproblems-problemphyxmini0283-lengthincentimeters}
  \lean{PhyXMiniProblems.ProblemPhyXMini0283.lengthInCentimeters}
  \uses{def:physics:phyx-mini-0283:phyxminiproblems-problemphyxmini0283-lengthquantity, def:physics:phyx-mini-0283:phyxminiproblems-problemphyxmini0283-quantityreadout, def:physics:phyx-mini-0283:phyxminiproblems-problemphyxmini0283-centimeterunitchoices}
  Centimeter readout of a physical length
\end{definition}
\begin{proof}
  This is the declared model definition of length In Centimeters.
\end{proof}
\begin{definition}[acceleration In Meters Per Second Squared]
  \label{def:physics:phyx-mini-0283:phyxminiproblems-problemphyxmini0283-accelerationinmeterspersecondsquared}
  \lean{PhyXMiniProblems.ProblemPhyXMini0283.accelerationInMetersPerSecondSquared}
  \uses{def:physics:phyx-mini-0283:phyxminiproblems-problemphyxmini0283-accelerationmagnitudequantity, def:physics:phyx-mini-0283:phyxminiproblems-problemphyxmini0283-quantityreadout}
  Meters-per-second-squared readout of an acceleration magnitude
\end{definition}
\begin{proof}
  This is the declared model definition of acceleration In Meters Per Second Squared.
\end{proof}
\begin{definition}[spring Stiffness In Newtons Per Meter]
  \label{def:physics:phyx-mini-0283:phyxminiproblems-problemphyxmini0283-springstiffnessinnewtonspermeter}
  \lean{PhyXMiniProblems.ProblemPhyXMini0283.springStiffnessInNewtonsPerMeter}
  \uses{def:physics:phyx-mini-0283:phyxminiproblems-problemphyxmini0283-springstiffnessquantity, def:physics:phyx-mini-0283:phyxminiproblems-problemphyxmini0283-quantityreadout}
  Newtons-per-meter readout of a spring stiffness
\end{definition}
\begin{proof}
  This is the declared model definition of spring Stiffness In Newtons Per Meter.
\end{proof}
\begin{definition}[Cord Side]
  \label{def:physics:phyx-mini-0283:phyxminiproblems-problemphyxmini0283-cordside}
  \lean{PhyXMiniProblems.ProblemPhyXMini0283.CordSide}
  The two modeled sides of the jump-seat suspension
\end{definition}
\begin{proof}
  This is the declared model definition of Cord Side.
\end{proof}
\begin{definition}[Figure Component]
  \label{def:physics:phyx-mini-0283:phyxminiproblems-problemphyxmini0283-figurecomponent}
  \lean{PhyXMiniProblems.ProblemPhyXMini0283.FigureComponent}
  Individually identifiable components visible in the supplied bitmap
\end{definition}
\begin{proof}
  This is the declared model definition of Figure Component.
\end{proof}
\begin{definition}[Frame Attachment]
  \label{def:physics:phyx-mini-0283:phyxminiproblems-problemphyxmini0283-frameattachment}
  \lean{PhyXMiniProblems.ProblemPhyXMini0283.FrameAttachment}
  Upper anchor positions on the horizontal support
\end{definition}
\begin{proof}
  This is the declared model definition of Frame Attachment.
\end{proof}
\begin{definition}[Seat Attachment]
  \label{def:physics:phyx-mini-0283:phyxminiproblems-problemphyxmini0283-seatattachment}
  \lean{PhyXMiniProblems.ProblemPhyXMini0283.SeatAttachment}
  Lower attachment positions on the seat
\end{definition}
\begin{proof}
  This is the declared model definition of Seat Attachment.
\end{proof}
\begin{definition}[Motion Axis]
  \label{def:physics:phyx-mini-0283:phyxminiproblems-problemphyxmini0283-motionaxis}
  \lean{PhyXMiniProblems.ProblemPhyXMini0283.MotionAxis}
  The direction in which the modeled seat is free to oscillate
\end{definition}
\begin{proof}
  This is the declared model definition of Motion Axis.
\end{proof}
\begin{definition}[Cord Mechanical Model]
  \label{def:physics:phyx-mini-0283:phyxminiproblems-problemphyxmini0283-cordmechanicalmodel}
  \lean{PhyXMiniProblems.ProblemPhyXMini0283.CordMechanicalModel}
  Mechanical idealization assigned to each elastic cord
\end{definition}
\begin{proof}
  This is the declared model definition of Cord Mechanical Model.
\end{proof}
\begin{definition}[Jump Seat Figure]
  \label{def:physics:phyx-mini-0283:phyxminiproblems-problemphyxmini0283-jumpseatfigure}
  \lean{PhyXMiniProblems.ProblemPhyXMini0283.JumpSeatFigure}
  \uses{def:physics:phyx-mini-0283:phyxminiproblems-problemphyxmini0283-cordside, def:physics:phyx-mini-0283:phyxminiproblems-problemphyxmini0283-figurecomponent, def:physics:phyx-mini-0283:phyxminiproblems-problemphyxmini0283-frameattachment, def:physics:phyx-mini-0283:phyxminiproblems-problemphyxmini0283-seatattachment}
  Qualitative and incidence data transcribed from the primary image. The image contains no printed labels or numerical scale; its colors and proportions do not supply any physical measurement
\end{definition}
\begin{proof}
  This is the declared model definition of Jump Seat Figure.
\end{proof}
\begin{definition}[Jump Seat Oscillator Setup]
  \label{def:physics:phyx-mini-0283:phyxminiproblems-problemphyxmini0283-jumpseatoscillatorsetup}
  \lean{PhyXMiniProblems.ProblemPhyXMini0283.JumpSeatOscillatorSetup}
  \uses{def:physics:phyx-mini-0283:phyxminiproblems-problemphyxmini0283-massquantity, def:physics:phyx-mini-0283:phyxminiproblems-problemphyxmini0283-lengthquantity, def:physics:phyx-mini-0283:phyxminiproblems-problemphyxmini0283-accelerationmagnitudequantity, def:physics:phyx-mini-0283:phyxminiproblems-problemphyxmini0283-springstiffnessquantity, def:physics:phyx-mini-0283:phyxminiproblems-problemphyxmini0283-cordside, def:physics:phyx-mini-0283:phyxminiproblems-problemphyxmini0283-motionaxis, def:physics:phyx-mini-0283:phyxminiproblems-problemphyxmini0283-cordmechanicalmodel, def:physics:phyx-mini-0283:phyxminiproblems-problemphyxmini0283-jumpseatfigure}
  All response quantities are independent fields. In particular, staticDescent is not defined from the answer. The scalar Physlib oscillator models the combined child-seat motion about the loaded equilibrium; later governing-law fields relate its positive mass and stiffness parameters to the dimensionful apparatus
\end{definition}
\begin{proof}
  This is the declared model definition of Jump Seat Oscillator Setup.
\end{proof}
\begin{definition}[Matches Problem And Primary Figure]
  \label{def:physics:phyx-mini-0283:phyxminiproblems-problemphyxmini0283-matchesproblemandprimaryfigure}
  \lean{PhyXMiniProblems.ProblemPhyXMini0283.MatchesProblemAndPrimaryFigure}
  \uses{def:physics:phyx-mini-0283:phyxminiproblems-problemphyxmini0283-lengthincentimeters, def:physics:phyx-mini-0283:phyxminiproblems-problemphyxmini0283-accelerationinmeterspersecondsquared, def:physics:phyx-mini-0283:phyxminiproblems-problemphyxmini0283-jumpseatoscillatorsetup}
  Problem data and primary-image readouts. The condition on maximum acceleration says that the design is evaluated at the stated 0.20 g safety threshold; it supplies no value for staticDescent
\end{definition}
\begin{proof}
  This is the declared model definition of Matches Problem And Primary Figure.
\end{proof}
\begin{definition}[Has Physical Jump Seat Parameters]
  \label{def:physics:phyx-mini-0283:phyxminiproblems-problemphyxmini0283-hasphysicaljumpseatparameters}
  \lean{PhyXMiniProblems.ProblemPhyXMini0283.HasPhysicalJumpSeatParameters}
  \uses{def:physics:phyx-mini-0283:phyxminiproblems-problemphyxmini0283-massinkilograms, def:physics:phyx-mini-0283:phyxminiproblems-problemphyxmini0283-lengthinmeters, def:physics:phyx-mini-0283:phyxminiproblems-problemphyxmini0283-accelerationinmeterspersecondsquared, def:physics:phyx-mini-0283:phyxminiproblems-problemphyxmini0283-springstiffnessinnewtonspermeter, def:physics:phyx-mini-0283:phyxminiproblems-problemphyxmini0283-jumpseatoscillatorsetup}
  Positivity and nondegeneracy conditions for the physical design
\end{definition}
\begin{proof}
  This is the declared model definition of Has Physical Jump Seat Parameters.
\end{proof}
\begin{definition}[Satisfies Two Cord Jump Seat Laws]
  \label{def:physics:phyx-mini-0283:phyxminiproblems-problemphyxmini0283-satisfiestwocordjumpseatlaws}
  \lean{PhyXMiniProblems.ProblemPhyXMini0283.SatisfiesTwoCordJumpSeatLaws}
  \uses{def:physics:phyx-mini-0283:phyxminiproblems-problemphyxmini0283-massinkilograms, def:physics:phyx-mini-0283:phyxminiproblems-problemphyxmini0283-lengthinmeters, def:physics:phyx-mini-0283:phyxminiproblems-problemphyxmini0283-accelerationinmeterspersecondsquared, def:physics:phyx-mini-0283:phyxminiproblems-problemphyxmini0283-springstiffnessinnewtonspermeter, def:physics:phyx-mini-0283:phyxminiproblems-problemphyxmini0283-jumpseatoscillatorsetup}
  The two Hookean cords undergo the same vertical extension and act in parallel, so their stiffnesses add. Static force balance gives K d\_s = m g. Around that loaded equilibrium, Physlib's harmonic oscillator has the same combined mass and effective stiffness, and an amplitude d\_m has maximum acceleration ω² d\_m. These are general mechanical laws. None states d\_s = 50 cm or chooses an answer label
\end{definition}
\begin{proof}
  This is the declared model definition of Satisfies Two Cord Jump Seat Laws.
\end{proof}
\begin{lemma}[maximum acceleration static descent relation]
  \label{lem:physics:phyx-mini-0283:phyxminiproblems-problemphyxmini0283-maximum-acceleration-static-descent-relation}
  \lean{PhyXMiniProblems.ProblemPhyXMini0283.maximum_acceleration_static_descent_relation}
  \uses{def:physics:phyx-mini-0283:phyxminiproblems-problemphyxmini0283-lengthinmeters, def:physics:phyx-mini-0283:phyxminiproblems-problemphyxmini0283-accelerationinmeterspersecondsquared, def:physics:phyx-mini-0283:phyxminiproblems-problemphyxmini0283-jumpseatoscillatorsetup, def:physics:phyx-mini-0283:phyxminiproblems-problemphyxmini0283-hasphysicaljumpseatparameters, def:physics:phyx-mini-0283:phyxminiproblems-problemphyxmini0283-satisfiestwocordjumpseatlaws}
  Eliminating the common mass and effective stiffness between static balance and the harmonic-oscillator acceleration law gives a\_max d\_s = g d\_m
\end{lemma}
\begin{proof}
  The conclusion follows from the governing relations and figure readouts named in the statement of maximum acceleration static descent relation.
\end{proof}
\begin{lemma}[static descent is fifty centimeters]
  \label{lem:physics:phyx-mini-0283:phyxminiproblems-problemphyxmini0283-static-descent-is-fifty-centimeters}
  \lean{PhyXMiniProblems.ProblemPhyXMini0283.static_descent_is_fifty_centimeters}
  \uses{def:physics:phyx-mini-0283:phyxminiproblems-problemphyxmini0283-lengthincentimeters, def:physics:phyx-mini-0283:phyxminiproblems-problemphyxmini0283-jumpseatoscillatorsetup, def:physics:phyx-mini-0283:phyxminiproblems-problemphyxmini0283-matchesproblemandprimaryfigure, def:physics:phyx-mini-0283:phyxminiproblems-problemphyxmini0283-hasphysicaljumpseatparameters, def:physics:phyx-mini-0283:phyxminiproblems-problemphyxmini0283-satisfiestwocordjumpseatlaws}
  The stated 10 cm amplitude and 0.20 g limit require d\_s = 50 cm
\end{lemma}
\begin{proof}
  The conclusion follows from the governing relations and figure readouts named in the statement of static descent is fifty centimeters.
\end{proof}
\begin{definition}[Answer Choice]
  \label{def:physics:phyx-mini-0283:phyxminiproblems-problemphyxmini0283-answerchoice}
  \lean{PhyXMiniProblems.ProblemPhyXMini0283.AnswerChoice}
  Labels of the four answers printed in the exercise
\end{definition}
\begin{proof}
  This is the declared model definition of Answer Choice.
\end{proof}
\begin{definition}[static Descent Centimeters]
  \label{def:physics:phyx-mini-0283:phyxminiproblems-problemphyxmini0283-answerchoice-staticdescentcentimeters}
  \lean{PhyXMiniProblems.ProblemPhyXMini0283.AnswerChoice.staticDescentCentimeters}
  \uses{def:physics:phyx-mini-0283:phyxminiproblems-problemphyxmini0283-answerchoice}
  Static-descent readout printed beside each answer, in centimeters
\end{definition}
\begin{proof}
  This is the declared model definition of static Descent Centimeters.
\end{proof}
\begin{definition}[recorded Answer Choice]
  \label{def:physics:phyx-mini-0283:phyxminiproblems-problemphyxmini0283-recordedanswerchoice}
  \lean{PhyXMiniProblems.ProblemPhyXMini0283.recordedAnswerChoice}
  \uses{def:physics:phyx-mini-0283:phyxminiproblems-problemphyxmini0283-answerchoice}
  The answer label recorded by the supplied dataset metadata
\end{definition}
\begin{proof}
  This is the declared model definition of recorded Answer Choice.
\end{proof}
\begin{definition}[Matches Answer Choice]
  \label{def:physics:phyx-mini-0283:phyxminiproblems-problemphyxmini0283-matchesanswerchoice}
  \lean{PhyXMiniProblems.ProblemPhyXMini0283.MatchesAnswerChoice}
  \uses{def:physics:phyx-mini-0283:phyxminiproblems-problemphyxmini0283-lengthincentimeters, def:physics:phyx-mini-0283:phyxminiproblems-problemphyxmini0283-jumpseatoscillatorsetup, def:physics:phyx-mini-0283:phyxminiproblems-problemphyxmini0283-answerchoice}
  A modeled design agrees with one of the displayed centimeter choices
\end{definition}
\begin{proof}
  This is the declared model definition of Matches Answer Choice.
\end{proof}
% --- Archon declaration topology end ---
% --- Archon physics formalization source end ---
